\section{Universal hyperplanes and the missing simultaneous Bertini producer}

Fix the original closed projective embedding $i:X\hookrightarrow\mathbf P^m_k$.
Write $L=i^*\mathcal O(1)$ and let $s_0,\ldots,s_m$ be its actual pulled
coordinate sections.  The affine coefficient parameter scheme is
$V=\operatorname{Spec}k[t_0,\ldots,t_m]$.  On the original product
$X_V=X\times_k V$, form the actual section
$\sigma=\sum_{a=0}^m t_a\,p_X^*s_a$ of $p_X^*L$.
Its actual zero-ideal subscheme $\mathcal I\subset X_V$ is the universal
hyperplane incidence, with its two original product projections.
The zero coefficient is retained; projectivizing the nonzero coefficient
locus is a subsequent quotient, not part of this affine construction.

\begin{definition}[The actual original universal incidence]
\label{def:kn-nadel-bertini-universal-incidence}
\lean{EnvelopeProperty.Vanishing.NadelBertiniUniversal.parameterRing,
EnvelopeProperty.Vanishing.NadelBertiniUniversal.product,
EnvelopeProperty.Vanishing.NadelBertiniUniversal.productLine,
EnvelopeProperty.Vanishing.NadelBertiniUniversal.coefficient,
EnvelopeProperty.Vanishing.NadelBertiniUniversal.universalSection,
EnvelopeProperty.Vanishing.NadelBertiniUniversal.incidenceIdeal,
EnvelopeProperty.Vanishing.NadelBertiniUniversal.incidence,
EnvelopeProperty.Vanishing.NadelBertiniUniversal.toSource,
EnvelopeProperty.Vanishing.NadelBertiniUniversal.toParameters,
EnvelopeProperty.Vanishing.NadelBertiniUniversal.projections_commute,
EnvelopeProperty.Vanishing.NadelBertiniUniversal.incidence_closed,
EnvelopeProperty.Vanishing.NadelBertiniUniversal.projectiveLine,
EnvelopeProperty.Vanishing.NadelBertiniUniversal.projectiveSections,
EnvelopeProperty.Vanishing.NadelBertiniUniversal.projectiveIncidence,
EnvelopeProperty.Vanishing.NadelBertiniUniversal.projectiveToSource,
EnvelopeProperty.Vanishing.NadelBertiniUniversal.projectiveToParameters}
\leanok
The original universal incidence and its two projections are constructed by
the displayed universal section and its actual canonical zero ideal.
The inclusion into the original scalar product is a closed immersion, and
the two original projections commute over the original coefficient base.
\end{definition}
\begin{proof}
Pull back each original native coordinate section to the same original
source and then to its actual coefficient product. Form their actual
linear combination with the original polynomial coefficient variables.
The actual image of dual evaluation is its genuine quasicoherent zero
ideal. Its actual ideal-sheaf subscheme is the required incidence. The
subscheme inclusion is closed, and its two composites with the genuine
product projections satisfy the original product square.
\end{proof}

\begin{lemma}[Eliminating the actual pivot coefficient]
\label{lem:kn-nadel-bertini-incidence-algebra}
\lean{EnvelopeProperty.Vanishing.NadelBertiniIncidenceAlgebra.remainder,
EnvelopeProperty.Vanishing.NadelBertiniIncidenceAlgebra.chartEquation,
EnvelopeProperty.Vanishing.NadelBertiniIncidenceAlgebra.chartIdeal,
EnvelopeProperty.Vanishing.NadelBertiniIncidenceAlgebra.chartEquiv,
EnvelopeProperty.Vanishing.NadelBertiniIncidenceAlgebra.chart_smooth,
EnvelopeProperty.Vanishing.NadelBertiniIncidenceAlgebra.chartProjection_smooth}
\leanok
Let $B$ be any commutative ring and let $(c_a)_{a\in S}$ be any finite
family of actual elements of $B$.  The actual quotient
\[
 B[t_*,(t_a)_{a\in S}]/\bigl(t_*+\sum_{a\in S}c_a t_a\bigr)
 \quad\simeq_B\quad B[(t_a)_{a\in S}]
\]
is obtained by retaining every original nonpivot variable and substituting
$t_*=-\sum c_a t_a$.  Its actual structure morphism to
$\operatorname{Spec}B$ is smooth.
\end{lemma}
\begin{proof}
The actual multivariable polynomial equivalence which isolates the pivot
variable carries the displayed relation to the monic linear polynomial
$T-(-\sum c_a t_a)$ over $B[(t_a)]$.  The induced quotient algebra
isomorphism followed by actual evaluation at $-\sum c_a t_a$ gives the
claimed equivalence, including the original $B$ structure.  Polynomial
algebras in finitely many variables are smooth over $B$. Transport
smoothness by this genuine algebra equivalence and then apply the actual
ring-to-affine-scheme smoothness comparison.
\end{proof}

\begin{lemma}[The actual coordinate sections give their own original frames]
\label{lem:kn-nadel-bertini-projective-coordinate-frame}
\lean{EnvelopeProperty.Vanishing.NadelBertiniProjectiveFrames.sourceOpen,
EnvelopeProperty.Vanishing.NadelBertiniProjectiveFrames.sectionHom,
EnvelopeProperty.Vanishing.NadelBertiniProjectiveFrames.section_isIso,
EnvelopeProperty.Vanishing.NadelBertiniProjectiveFrames.coordinateFrame,
EnvelopeProperty.Vanishing.NadelBertiniProjectiveFrames.coordinateFrame_section,
EnvelopeProperty.Vanishing.NadelBertiniProjectiveFrames.sourceOpen_affine}
\leanok
\uses{def:kn-nadel-bertini-universal-incidence}
On the actual original coordinate preimage $U_j$, the actual pulled
coordinate section is an isomorphism from the original structure line to
$L|_{U_j}$. Its actual inverse is a frame in which that same coordinate
section is exactly one. For a closed embedding $i$, $U_j$ is affine.
\end{lemma}
\begin{proof}
The native $j$th coordinate section is an isomorphism on its own original
ambient chart. The genuine restriction/pullback comparison retains that
whole map. Pull it back along the actual restriction of $i$ to $U_j$;
whole-section pullback composition identifies the result with the same
original coordinate section on $U_j$. Its inverse gives the claimed actual
frame and the inverse identity gives the exact normalization. A closed
immersion is affine and the ambient coordinate chart is affine, proving
that its original preimage is affine.
\end{proof}

\begin{lemma}[The actual original affine coefficient product ring]
\label{lem:kn-nadel-bertini-polynomial-product}
\lean{EnvelopeProperty.Vanishing.NadelBertiniPolynomial.sectionsEquiv,
EnvelopeProperty.Vanishing.NadelBertiniPolynomial.sectionsEquiv_X,
EnvelopeProperty.Vanishing.NadelBertiniPolynomial.sectionsEquiv_C,
EnvelopeProperty.Vanishing.NadelBertiniPolynomial.sectionsEquiv_eval,
EnvelopeProperty.Vanishing.NadelBertiniPolynomial.sectionsEquiv_linearForm}
\leanok
\uses{def:kn-nadel-bertini-universal-incidence}
For an actual original affine scheme $U$ over $\operatorname{Spec}R$, the
whole ring of the actual original product $U\times_R\mathbf A^S_R$
is the original polynomial ring $\Gamma(U)[(t_a)]$. Its actual isomorphism
sends each original variable to the second-projection coefficient and each
original constant to the first-projection section. It therefore retains
the complete original universal linear polynomial.
\end{lemma}
\begin{proof}
The original affine product's two actual projections construct the genuine
pushout square of rings of sections, without a flatness premise. The
original tensor-product pushout therefore identifies its whole section
ring with $R[(t_a)]\otimes_R\Gamma(U)$. Commute the actual two factors
and apply the genuine polynomial-algebra tensor equivalence. Its actual
inverse retains the original variable as $1\otimes t_a$ and each original
constant as $b\otimes1$. The original affine product comparison then gives
the two stated projection formulas. Polynomial ring-map extensionality
gives the whole evaluation formula, and actual sum and product preservation
gives the entire linear-polynomial identity.
\end{proof}

\begin{lemma}[The actual universal section retains its original frame coefficients]
\label{lem:kn-nadel-bertini-universal-frame-coefficient}
\lean{EnvelopeProperty.Vanishing.NadelBertiniFrameCoefficients.unit_pullback_scalar,
EnvelopeProperty.Vanishing.NadelBertiniFrameCoefficients.pulledFrame,
EnvelopeProperty.Vanishing.NadelBertiniFrameCoefficients.pulledFrame_section,
EnvelopeProperty.Vanishing.NadelBertiniFrameCoefficients.universal_frame_coefficient}
\leanok
\uses{def:kn-nadel-bertini-universal-incidence,lem:kn-nadel-bertini-polynomial-product,lem:kn-nadel-bertini-projective-coordinate-frame}
Let $U$ be an actual original affine source carrying an actual line frame.
The actual pulled frame on its original coefficient product sends the
whole original universal section to the image of
$\sum_a c_a t_a$ under the preceding actual product-ring equivalence,
where $c_a$ is the coefficient of the same original section in that frame.
\end{lemma}
\begin{proof}
The native unit-section pullback commutes with the original whole frame
map. Its original structure-line unit comparison sends a pulled scalar
section to its actual structure-ring pullback. This follows by pulling
the whole structure-line scalar endomorphism and evaluating it at one.
Thus the actual pulled frame sends each original pulled section to the
actual pullback of its original coefficient. It retains sums and scalar
multiplication. Apply that identity to every term of the actual universal
section, commute the two original scalar factors, and use the complete
linear-polynomial formula for the actual affine coefficient product.
\end{proof}

\begin{lemma}[The actual whole original affine incidence ideal]
\label{lem:kn-nadel-bertini-affine-incidence-ideal}
\lean{EnvelopeProperty.Vanishing.NadelBertiniAffineIncidence.zeroIdeal_top,
EnvelopeProperty.Vanishing.NadelBertiniAffineIncidence.incidenceIdeal_top}
\leanok
\uses{def:kn-nadel-bertini-universal-incidence,lem:kn-nadel-bertini-universal-frame-coefficient}
On an actual original affine source with its actual line frame, the whole
original universal incidence ideal on the full coefficient product is the
image of the actual principal ideal $(\sum c_a t_a)$ under the proved
original polynomial-product ring equivalence.
\end{lemma}
\begin{proof}
The original line-section zero ideal on an affine frame is precisely the
principal ideal of the actual frame coefficient. Apply this genuine whole
ideal statement on the original coefficient product, with the actual
pulled frame and the actual universal section. The preceding exact frame
coefficient theorem gives the image of the entire original linear
polynomial. The image of its actual principal ideal is the principal ideal
of its image. This identifies the whole original ideal, retaining all
scheme structure and without passing to its radical.
\end{proof}

\begin{lemma}[Any actual unit pivot makes the original linear quotient smooth]
\label{lem:kn-nadel-bertini-linear-pivot-smooth}
\lean{EnvelopeProperty.Vanishing.NadelBertiniLinearSmooth.linearForm,
EnvelopeProperty.Vanishing.NadelBertiniLinearSmooth.linearIdeal,
EnvelopeProperty.Vanishing.NadelBertiniLinearSmooth.linearForm_pderiv,
EnvelopeProperty.Vanishing.NadelBertiniLinearSmooth.linear_smooth}
\leanok
For a commutative ring $B$ and finite small variable set $S$, let
$F=\sum_{a\in S}c_at_a$. If the actual original coefficient $c_j$ is one,
the original structure algebra $B\to B[(t_a)]/(F)$ is smooth.
\end{lemma}
\begin{proof}
Take the genuine polynomial presentation with the single original relation
$F$. Assign its unique relation to the actual original pivot variable $j$.
Its original partial derivative is $C(c_j)=1$. The resulting actual one by
one Jacobian matrix is the identity, so its Jacobian is a unit in the
original quotient. This constructs an actual submersive presentation of
the original quotient, proving standard smoothness and then smoothness.
No quotient is replaced by its reduction and no smoothness certificate is
supplied.
\end{proof}

\begin{lemma}[Actual affine subscheme sections retain the original projection algebra]
\label{lem:kn-nadel-bertini-affine-subscheme-sections}
\lean{EnvelopeProperty.Vanishing.NadelBertiniSubschemeSections.quotientEquivSections,
EnvelopeProperty.Vanishing.NadelBertiniSubschemeSections.quotientEquivSections_mk,
EnvelopeProperty.Vanishing.NadelBertiniSubschemeSections.quotientAlgEquivSections}
\leanok
For any original affine scheme $Y$, whole original ideal sheaf $I$, and
original morphism $g:Y\to X$, the actual closed subscheme's original
section ring is isomorphic to $\Gamma(Y)/I(\top)$. It carries each original
quotient class of $b$ to its actual restricted section. This constructs an
actual algebra equivalence over $\Gamma(X)$ for the same original composite
$I\hookrightarrow Y\to X$.
\end{lemma}
\begin{proof}
The actual subscheme inclusion is a closed immersion with affine target,
so it is surjective on global sections. The original whole ideal sheaf is
its actual kernel; on an affine scheme this is exactly the kernel of the
actual global-section map. The genuine quotient-kernel first isomorphism
gives the ring equivalence and its entire original class formula. The
original projection's composite formula on global sections proves
compatibility with each original $\Gamma(X)$ scalar. This constructs the
actual algebra equivalence rather than assuming such compatibility.
\end{proof}

\begin{lemma}[Actual original affine incidence projection is smooth]
\label{lem:kn-nadel-bertini-affine-incidence-smooth}
\lean{EnvelopeProperty.Vanishing.NadelBertiniAffineSmooth.sourceSectionsAlgEquiv,
EnvelopeProperty.Vanishing.NadelBertiniAffineSmooth.affine_toSource_smooth}
\leanok
\uses{lem:kn-nadel-bertini-affine-incidence-ideal,lem:kn-nadel-bertini-linear-pivot-smooth,lem:kn-nadel-bertini-affine-subscheme-sections}
For an actual affine original source $U$, original line $L$ with an actual
whole frame, and original section family $(s_a)$, suppose the coefficient
of one original section in that frame is exactly one. The projection of
the actual whole universal incidence to this same $U$ is smooth.
\end{lemma}
\begin{proof}
Upgrade the original polynomial-product ring equivalence to an algebra
equivalence over the same original $\Gamma(U)$, using its already proved
constant-section formula. The entire original incidence ideal is the image
of the whole original linear ideal. The induced quotient algebra
equivalence transports the proved actual pivot smoothness to the original
product-ring quotient. The actual affine subscheme-section algebra
equivalence then transports it to the section ring of the same original
incidence. The actual subscheme is affine because its original inclusion
is closed with affine target. The affine ring-homomorphism criterion for
smoothness now applies to its original source projection, whose scalar
map is precisely the map retained throughout both algebra equivalences.
\end{proof}

\begin{lemma}[Native whole-section pullback retains its actual linear combination]
\label{lem:kn-nadel-bertini-pullback-sections}
\lean{EnvelopeProperty.Vanishing.NadelBertiniPullbackSections.pullback_add,
EnvelopeProperty.Vanishing.NadelBertiniPullbackSections.pullback_smul,
EnvelopeProperty.Vanishing.NadelBertiniPullbackSections.frameViaFactor_unitSection}
\leanok
\uses{lem:kn-nadel-bertini-universal-frame-coefficient}
For every actual scheme morphism, the native whole-section pullback
preserves addition and carries $r s$ to the actual pulled scalar times the
actual pulled section. If the same morphism factors through an original
local frame, its actual factor frame sends the native pulled section to
the actual pullback of the original local frame coefficient.
\end{lemma}
\begin{proof}
The genuine pullback-adjunction unit is linear on the original module,
where its pushforward scalar action is the original structure-ring
pullback. The final transport from the inverse image of the whole open to
the whole open preserves sums and transports exactly the same scalar.
For the last assertion, replace the original whole section by its actual
structure-line section morphism. The native factor-frame pairing carries
this same original section morphism to the pulled structure-line scalar
endomorphism. Evaluating the actual identity at one and using the original
pullback-unit section comparison gives the claimed coefficient formula.
\end{proof}

\begin{lemma}[The actual original incidence is smooth over its original source]
\label{lem:kn-nadel-bertini-incidence-bundle}
\uses{def:kn-nadel-bertini-universal-incidence,lem:kn-nadel-bertini-incidence-algebra,lem:kn-nadel-bertini-projective-coordinate-frame,lem:kn-nadel-bertini-polynomial-product,lem:kn-nadel-bertini-universal-frame-coefficient,lem:kn-nadel-bertini-affine-incidence-ideal,lem:kn-nadel-bertini-linear-pivot-smooth,lem:kn-nadel-bertini-affine-subscheme-sections,lem:kn-nadel-bertini-affine-incidence-smooth,lem:kn-nadel-bertini-pullback-sections}
The projection of the actual original universal incidence $\mathcal I$
to the same original source $X$ is smooth. If $X$ is smooth over $k$,
then so is the same original $\mathcal I$ over $k$.
\end{lemma}
\begin{proof}
The original coordinate-section trivialization and actual pullback-unit
naturality identify the coefficient of $\sigma$ on each affine subopen
of $U_j$ with the displayed polynomial.  The actual line-section zero
ideal is the principal ideal of this coefficient.  The preceding actual
quotient equivalence identifies its full subscheme chart with affine
$m$-space over that original affine subopen.  The $U_j$ cover $X$ and
their inverse images cover the incidence.  Smoothness is local on source
and target, so these actual identifications prove the assertion globally.
Composition with the original smooth structure map gives the last claim.
\end{proof}

\begin{lemma}[A genuine open of smooth original hyperplane sections]
\label{lem:kn-nadel-bertini-smooth-parameter-open}
\uses{lem:kn-nadel-bertini-incidence-bundle}
Suppose $k$ has characteristic zero and $X$ is smooth and projective over
$k$.  There is a nonempty open in the original coefficient parameter
scheme on which every original universal hyperplane fibre is smooth.
The same conclusion holds simultaneously for any finite family of smooth
closed projective subschemes, using their actual restricted sections.
\end{lemma}
\begin{proof}
The original incidence is smooth over $k$ by the preceding lemma. Over the
generic point of the parameter affine space, its affine coordinate rings
are localizations of regular rings. They remain regular, and are smooth
over the characteristic-zero generic field. Thus the actual nonsmooth
locus of the parameter projection misses the generic fibre. That locus is
closed; the parameter projection is proper because the original incidence
is closed in $X\times V$ and $X$ is projective. Its image is therefore
closed and misses the original generic point. The complement is the
required nonempty original open. Apply the same argument to each original
closed subscheme and intersect the finitely many resulting opens. The
original parameter affine space is integral, so its generic point belongs
to every such open. This constructs their common nonempty open.
\end{proof}

\begin{theorem}[The required simultaneous auxiliary divisors]
\label{thm:kn-nadel-bertini-auxiliary-divisors}
\uses{lem:kn-nadel-bertini-smooth-parameter-open}
Let $X$ be the original smooth projective $d$-dimensional scheme over a
characteristic-zero field, with the original finite
SNC boundary $(D_i)$. For the original ample line bundles $A_i$, all
sufficiently large common exponents $N$ allow actual divisors
$H_{i,j}\in|A_i^N(-D_i)|$, $0\le j\le d$, such that the total original
boundary remains SNC and, for each $i$, the original common intersection
$\bigcap_{j=0}^d H_{i,j}$ is empty.
\end{theorem}
\begin{proof}
For each of the finitely many $i$, ampleness makes $A_i^N(-D_i)$ very
ample for all sufficiently large $N$; take one common bound. At each
finite step use its actual closed projective embedding. Apply the previous
parameter-open construction to all closed strata of the current original
SNC boundary, including the intersections with the already selected
auxiliary divisors. Require also that no positive-dimensional stratum is
contained in the new hyperplane, and avoid the finitely many
zero-dimensional strata. These are open conditions containing the
parameter generic point. The infinite original field supplies an actual rational point in their
common nonempty affine parameter open: take a nonempty principal subopen
and use that a nonzero polynomial cannot vanish on all tuples over an
infinite field. Pulling the
universal section back to this actual point selects the required original
divisor. The differential transversality to every original stratum gives
actual SNC charts for the enlarged boundary. Each successive member for a
fixed $i$ lowers the dimension of its previous nonempty intersection by
one. After $d+1$ such members that intersection is empty. Carry out the
finite induction over all original indices, retaining the earlier selected
divisors at each step.
\end{proof}
